

\section{引言}
\subsection{问题背景}

零售行业作为国民经济的重要组成部分，其销售额的波动与变化直接反映了消费市场的景气程度。准确预测零售销售额对于企业经营决策、供应链管理、营销规划等方面具有至关重要的意义。随着大数据技术的快速发展和企业信息化水平的持续提升，零售企业积累了海量的交易数据，如何从这些数据中挖掘有价值的商业洞察成为学术界和产业界共同关注的核心议题。

零售销售额的变化受到多重因素的复杂影响，包括时间周期特征、产品品类结构、区域经济差异、客户消费行为以及物流履约效率等。传统上，零售企业依赖经验判断和简单的移动平均方法进行销售预测，精度有限且难以量化各因素的实际影响权重。近年来，时间序列分析和机器学习方法在经济预测和商业智能领域展现出巨大潜力，为零售销售额的精准预测提供了新的技术路径。

在此背景下，本文基于2022-2025年美国零售企业的全量订单数据集（9800条交易记录，18个核心字段），综合运用描述性统计、方差分析、相关性分析和时间序列预测等方法，系统研究零售销售额的数据特征、影响因素与变化规律，构建科学的预测模型并给出可落地的经营策略建议。

\subsection{问题重述}

\textbf{问题一：}对零售数据集进行基本统计与可视化分析，计算销售额的核心统计量及业务指标，绘制时间趋势图和多维度分布图。

\textbf{问题二：}从时间、地理、产品、客户、物流五个维度量化分析各因素与销售额的关联程度，绘制相关性可视化图表，识别驱动销售额增长的核心因素。

\textbf{问题三：}选择合适的建模方法构建零售销售额预测模型，按时间维度划分训练集与测试集，完成模型训练与参数优化并全面评估预测效果。

\textbf{问题四：}使用优化完成的预测模型对未来一段时间（2026年全年及分季度）的零售销售额进行预测，并结合影响因素分析结论，提供可落地的经营策略建议。
